\chapter{Life at FlowGenX AI}
\label{ch:life}

This chapter reflects on what it was actually like to work at
FlowGenX AI as an intern. Because my internship was fully remote,
this is necessarily a different kind of chapter from the
``life-at-the-office'' narratives common in earlier internship
reports. Remote-first work is now the dominant model in modern
software companies, and the lived experience of being an intern at
a globally distributed startup deserves to be described on its own
terms rather than as a poor substitute for an in-office one.

\section{A Remote-First, Globally Distributed Workplace}
FlowGenX AI is headquartered in San Francisco, California, with
engineers based in the United States and Bangladesh. There is no
shared physical office in Bangladesh; the company runs as a
remote-first organisation, with all employees and interns working
from their own homes or chosen workspaces. The expectations and
support structures the company puts in place are designed for that
reality: the tools, communication norms, and cultural cues
collectively make distributed work feel coherent rather than
fragmented.

The environment in practice has the following defining traits:

\begin{itemize}[leftmargin=*]
    \item \textbf{Trust over presence.} Nobody monitors when an
    engineer is online; what matters is whether the work moves
    forward, the pull requests get reviewed, and the standup updates
    are posted. This puts the burden of self-direction on the
    engineer, but in exchange offers a kind of working life that is
    very hard to find in traditional offices.

    \item \textbf{Asynchronous-first communication.} Most coordination
    happens through Google Chat threads and ClickUp comments. Meetings
    are reserved for design discussions, code-walkthroughs, and
    onboarding sessions. The team has been thoughtful about not
    over-scheduling video calls.

    \item \textbf{Globally distributed culture.} Working alongside
    engineers and leadership in San Francisco from Dhaka is itself a
    professional experience worth highlighting. It teaches habits ---
    written clarity, asynchronous handoffs, time-zone empathy ---
    that translate directly into how senior software roles operate
    today.
\end{itemize}

\section{My Workspace and Equipment}
Per my offer letter, the internship was full-time, remote
(Bangladesh), and I used my personal laptop for all work. I set up
a dedicated workspace at home with:

\begin{itemize}[leftmargin=*]
    \item A reliable broadband connection with a backup mobile
    hotspot for resilience.
    \item Local installations of Python 3.11, Docker Desktop, the
    \texttt{uv} package manager, Visual Studio Code, Postman, Git,
    and the GitHub CLI.
    \item Local clones of the four FlowGenX repositories under a
    single parent directory so that editable installs cross-referenced
    each other cleanly.
    \item A local PostgreSQL instance for connector-metadata
    development and testing.
\end{itemize}

The remote-first model meant I owned the operational health of my
own development environment. When something broke (Docker on Windows
having a moment, a Python virtual environment getting out of sync,
a PostgreSQL upgrade requiring a re-import), I learned to diagnose
and fix it myself rather than escalating to an IT team. This is not
a small skill --- and it is one that classroom education rarely
teaches.

\section{Daily and Weekly Cadence}
\subsection{A Typical Working Day}
My typical working day ran from approximately
\textbf{9:00\,a.m.\ to 5:00\,p.m.\ Bangladesh Standard Time
(UTC+6)}, with flexibility to shift hours when a synchronous meeting
with the United States team required late-evening overlap. A typical
day looked roughly like this:

\begin{itemize}[leftmargin=*]
    \item \textbf{Morning (9:00 -- 11:00).} Catch up on overnight
    activity from the United States team --- Google Chat messages,
    pull-request comments, ClickUp updates. Resolve straightforward
    review feedback before starting deep work.
    \item \textbf{Late morning to early afternoon (11:00 -- 2:30).}
    Focused implementation work on whichever connector or fix was the
    day's primary deliverable.
    \item \textbf{Afternoon (2:30 -- 5:00).} Testing, integration test
    runs, sample-backfill work, and any outstanding documentation.
    \item \textbf{Evening (occasional).} Synchronous meetings or
    code-walkthroughs with the United States team during their
    morning, when an overlap was needed.
\end{itemize}

\subsection{Weekly Rhythm}
The week alternated between focused individual work and lightweight
synchronous touchpoints. Pull requests typically opened and merged
through the week as connectors were ready, rather than being
batched into a single end-of-week push, and this rolling cadence is
what made the team's velocity feel sustainable rather than spiky.

\section{Communication and Collaboration Style}
\subsection{Tools}
\begin{itemize}[leftmargin=*]
    \item \textbf{Google Chat} for asynchronous messaging --- both
    one-on-one and team channels.
    \item \textbf{Google Meet} for synchronous discussions, design
    reviews, and onboarding sessions.
    \item \textbf{ClickUp} for sprint and task tracking.
    \item \textbf{GitHub} for code review through pull-request
    comments.
\end{itemize}

The team migrated from \emph{Microsoft Teams} to \emph{Google Chat}
during my tenure, and I participated in the migration first-hand ---
a small but real experience of how a distributed team rolls out
infrastructure changes without disrupting ongoing work.

\subsection{Norms}
A few unwritten norms shaped how the team worked: questions go into
public team channels (rather than direct messages) so the answers
remain searchable; messages are written to be self-contained so the
reader does not need additional context to act; help requests show
what has already been tried and observed rather than just ``it does
not work''; and questions are phrased with timezone empathy so they
can be answered overnight and come back actionable in the morning.

\section{Onboarding Experience}
The onboarding period was deliberately short and project-driven.
There was no multi-week classroom-style training; instead I was
pointed at the project's internal documentation, walked through the
architecture by my supervisor, and
encouraged to start building real connectors as my learning
artefacts. By the end of the second week I had merged my first
half-dozen pull requests and had a complete mental model of the
codebase.

For me personally this approach worked very well. Building real
features from week two meant I was investing in the same code I
would still be working on six months later, and the depth of
understanding that comes from shipping is qualitatively different
from the kind that comes from reading.

\section{Mentorship and Learning Culture}
The single most positive aspect of life at FlowGenX AI was the
mentorship. My supervisor reviewed every pull request thoroughly,
explained the reasoning behind every suggestion, and never made me
feel that a question was too small. My team lead's design feedback
on larger refactors gave me a sense of how senior engineers think
about evolving a codebase. My teammates were generous with their
time when I needed help understanding a piece of the runtime that
sat outside my immediate work.

Beyond direct mentorship, the company maintains an active engineering
blog (\texttt{https://blog.flowgenx.ai/}) and developer documentation
(\texttt{https://docs.flowgenx.ai/}) that engineers are encouraged to
read and contribute to. The Generative AI space moves fast, and
working at FlowGenX meant being constantly nudged toward what was new
in the field --- new agent patterns, new providers, new protocols.

\section{Cross-Time-Zone Collaboration and Reflection}
Working from Bangladesh on a team partly based in San Francisco ---
roughly a twelve-hour offset, with only a few precious overlap hours
per day --- required intentional handoffs. I learned to leave
detailed end-of-day updates so the United States team could pick up
where I left off, to compose questions that did not require an
immediate reply (so the answer could come back overnight and be
actionable in the morning), and to schedule the small number of
synchronous meetings during the evening overlap window without
disrupting either side's focus time. By the end of the internship,
this discipline felt entirely natural and is one of the most
concretely transferable skills the experience gave me.

There is sometimes a quiet assumption in the academic conversation
about internships that ``real'' internships happen in offices. My
experience challenged that assumption. What I missed --- the casual
hallway conversations, the in-person whiteboard sessions --- I missed
honestly. But what I gained was substantial: the experience of
contributing to a serious product alongside experienced engineers in
San Francisco, California, with the agency and trust that comes from
being treated as a contributing member of the team from day one.
